\section{Extended model results}\label{si:results}

\subsection{Fast-land yield sweep and parameter sensitivity of the composition diagnostic}\label{si:sweep}

Table~\ref{tab:si_sweep} extends the main-text fast-land-yield sweep in Table~2(a) by reporting the balance-branch slope. The sign of $\dconv$ stays positive when $b_{G,c}$ triples. At $b_f=0.85$\,gha\,ha$^{-1}$, $\Ups(q^*)=0.058,\,0.066,\,0.074$\,gha\,ha$^{-1}$ for $b_{G,c}=0.4,\,0.8,\,1.2$\,gha\,yr\,ha$^{-1}$, giving $\dconv=+0.792,\,+0.784,\,+0.776$\,gha\,ha$^{-1}$.

\begin{table}[htbp]
\centering\small
\caption{Fast-land yield sweep at the baseline (all other parameters fixed): composition diagnostic at $q^*$, eigenvalues of the deficit-side matrix, slope of the balance branch and demand ceiling.}
\label{tab:si_sweep}
\begin{tabular}{l r r r r r}
\toprule
$b_f$ (gha\,ha$^{-1}$) & $\dconv(q^*)$ (gha\,ha$^{-1}$) & $\lambda_1$ (yr$^{-1}$) & $\lambda_2$ (yr$^{-1}$) & $b_fk-b_c$ (gha\,ha$^{-1}$) & $\Eceil$ (gha) \\
\midrule
$0.60$ & $+0.534$ & $-0.0522$ & $-0.975$ & $0.277$ & $0.804$ \\
$0.85$ & $+0.784$ & $-0.0537$ & $-0.998$ & $0.414$ & $1.138$ \\
$1.20$ & $+1.134$ & $-0.0547$ & $-1.014$ & $0.605$ & $1.606$ \\
$2.00$ & $+1.934$ & $-0.0556$ & $-1.029$ & $1.041$ & $2.675$ \\
\bottomrule
\end{tabular}
\end{table}

\subsection{Drawdown-depth ladder}\label{si:ladder}

Table~\ref{tab:si_drawdown} reports the minimum capital-land area reached in endogenous-demand runs over a grid of initial populations $P_0$. The runs use $A_{f0}=0.6$, $q_0=0.9$, $\tau_g=20$, $\tau_p=25$\,yr, $\chi_r=0.10$, $\Delta t=0.1$\,yr, and $T=400$\,yr; the full $P_0$ grid (step $0.1$) is in \texttt{si\_drawdown\_depth.csv}. Depending on $P_0$, the minimum occurs 3--25\,yr after the run begins, when the shortfall first closes ($B\ge E$). At $A_{c0}=1.00$\,ha, the delayed-quality $10^{-6}$\,ha proximity values shift by at most $0.001$\,cap across the tested $\tau_g=5,20,40$\,yr cases (Table~\ref{tab:si_delay}). For the baseline history, the $\tau_g=20$ and $40$\,yr quality paths are exactly identical because both reach the imposed cap before either delay elapses; this is saturation-induced degeneracy, not general delay robustness. The $\tau_g=5$\,yr case is distinct. Drawdown depth varies continuously and monotonically with both $P_0$ and $A_{c0}$.

The initial populations at which the minimum comes within $10^{-2}$, $10^{-3}$, and $10^{-4}$\,ha of the floor are $1.8,\,1.9,\,2.1$\,cap for $A_{c0}=0.40$\,ha; $2.2,\,2.3,\,2.4$ for $0.70$\,ha; $2.6,\,2.7,\,2.8$ for $1.00$\,ha; and $2.8,\,2.9,\,3.0$ for $1.18$\,ha. Bisection with a $10^{-6}$\,ha proximity tolerance ($\Delta t=0.05$\,yr, $T=1000$\,yr) gives $2.41,\,2.70,\,2.98,\,3.14$\,cap. The thresholds match to three decimals for $\chi_r=0.10$ and $\chi_r=0$ ($2.387$ and $2.955$\,cap at $A_{c0}=0.40$ and $1.00$\,ha, $\Delta t=0.1$\,yr). \par\pagebreak

\begin{table}[htbp]
\centering\small
\caption{Minimum capital land $\min_tA_c$ (ha) reached during the run, by initial capital $A_{c0}$ (rows, ha) and initial population $P_0$ (columns, cap); typed floor $\Acmin=0.05$\,ha.}
\label{tab:si_drawdown}
\begin{tabular}{lrrrrrrrr}
\toprule
$A_{c0}$ $\backslash$ $P_0$ & 1.0 & 1.5 & 2.0 & 2.5 & 2.8 & 3.0 & 3.5 & 4.0 \\
\midrule
$0.4$ & 0.396 & 0.147 & 0.050 & 0.050 & 0.050 & 0.050 & 0.050 & 0.050 \\
$0.7$ & 0.700 & 0.442 & 0.121 & 0.050 & 0.050 & 0.050 & 0.050 & 0.050 \\
$1.0$ & 1.000 & 0.731 & 0.398 & 0.075 & 0.050 & 0.050 & 0.050 & 0.050 \\
$1.18$ & 1.180 & 0.898 & 0.557 & 0.230 & 0.055 & 0.050 & 0.050 & 0.050 \\
\bottomrule
\end{tabular}

\end{table}

\subsection{Recovery metric and the role of the restoration flow}\label{si:recovery}

We measure recovery as $\mathcal R_c(T)=(A_c(T)-A_c^{\mathrm{deg}})/(A_c^{\mathrm{init}}-A_c^{\mathrm{deg}})$. A value of 0 means no recovery, 1 means return to the reference, a negative value means further decline, and a value above 1 means over-restoration. Table~\ref{tab:si_recovery} applies this metric to delayed-quality runs with $A_{c0}=1.0$\,ha ($A_c^{\mathrm{init}}=1.0$\,ha, $A_c^{\mathrm{deg}}=\min_tA_c$, $T=400$\,yr). Without restoration, $A_c$ stays at its minimum because no return flow exists. With surplus-gated restoration, recovery is partial in all six initial-population cases ($\mathcal R_c=0.21$--$0.89$); none returns to the reference. The separate gate-sign experiment in SI~\S S4 uses other initial states and interventions, including cases of over-restoration.

\begin{table}[htbp]
\centering\small
\caption{Recovery of the capital-land area after a demand-driven drawdown ($A_{c0}=1.0$\,ha, $A_{f0}=0.6$\,ha, $q_0=0.9$, $T=400$\,yr): initial population $P_0$ is in cap, initial footprint $E_0$ in gha, capital-land values in ha, and the recovery metric is dimensionless.}
\label{tab:si_recovery}
\begin{tabular}{l r r r r r r}
\toprule
$P_0$ & $E_0$ & $\min_tA_c$ & $t(\min)$ (yr) & $A_c(T)$, $\chi_r=0$ & $A_c(T)$, $\chi_r=0.10$ & $\mathcal R_c(T)$, $\chi_r=0.10$ \\
\midrule
$1.5$ & $0.825$ & $0.731$ & $3.0$ & $0.731$ & $0.969$ & $0.89$ \\
$2.0$ & $1.100$ & $0.398$ & $3.8$ & $0.398$ & $0.677$ & $0.46$ \\
$2.5$ & $1.375$ & $0.075$ & $4.1$ & $0.075$ & $0.433$ & $0.39$ \\
$3.0$ & $1.650$ & $0.050$ & $13.8$ & $0.050$ & $0.328$ & $0.29$ \\
$3.5$ & $1.925$ & $0.050$ & $20.6$ & $0.050$ & $0.278$ & $0.24$ \\
$4.0$ & $2.200$ & $0.050$ & $26.6$ & $0.050$ & $0.250$ & $0.21$ \\
\bottomrule
\end{tabular}

\end{table}

\subsection{Gate-sign phase dependence}\label{si:gates}

The gate-sign comparison starts from a post-drawdown state: $A_c=0.058$\,ha (near the floor), $A_f=0.6$\,ha, $q=0.9$, and $P=1.2$\,cap. At the degraded equilibrium, $R_B=1$ and $B=E$, so neither restoration gate is active. The reference equilibrium has $A_c^{\mathrm{init}}=0.8225$\,ha and $P^{\mathrm{init}}=1.489$\,cap; the degraded reference is $A_c^{\mathrm{deg}}=0.0513$\,ha. Every row uses $T_b=0$ and the same starting state.

\begin{table}[H]
\centering\scriptsize
\setlength{\tabcolsep}{2pt}
\caption{Restoration gate and recovery outcomes at $T=100$ and $900$\,yr. The table reports end-state land areas, population relative to the reference, trajectory extrema, and peak debt $D$ (gha\,yr). Annual account values $B=E$ are in gha; land areas are in ha; population ratios and recovery metrics are dimensionless. Each row approaches $R_B=1$ ($E=B$).}
\label{tab:si_gates}
\begin{tabular}{l r r r r r r r r r r}
\toprule
Policy / gate & $\mathcal R_c(100)$ & $\mathcal R_c(900)$ & $A_c$ & $A_f$ & $A_r$ & $B=E$ & $P/P^{\mathrm{init}}$ & $\min A_r$ & $\max A_c$ & $D_{peak}$ \\
\midrule
none ($R_{rc}=0$) & $-0.00$ & $-0.00$ & $0.051$ & $1.336$ & $1.113$ & $1.138$ & $1.39$ & $1.113$ & $0.057$ & $0.152$ \\
surplus-gated, $\chi_r=0.02$ & $0.69$ & $0.73$ & $0.612$ & $1.030$ & $0.858$ & $0.906$ & $1.11$ & $0.858$ & $0.612$ & $0.152$ \\
surplus-gated, $\chi_r=0.05$ & $1.08$ & $1.07$ & $0.875$ & $0.886$ & $0.738$ & $0.797$ & $0.97$ & $0.724$ & $0.885$ & $0.152$ \\
surplus-gated, $\chi_r=0.10$ & $1.28$ & $1.27$ & $1.030$ & $0.802$ & $0.668$ & $0.733$ & $0.90$ & $0.602$ & $1.092$ & $0.152$ \\
surplus, $\chi_r=0.10$, target cap $A_c^*=0.8225$ & $1.00$ & $1.00$ & $0.822$ & $0.915$ & $0.763$ & $0.819$ & $1.00$ & $0.664$ & $0.822$ & $0.152$ \\
surplus, $\chi_r=0.10$, conversion frozen & $1.35$ & $1.36$ & $1.097$ & $0.765$ & $0.638$ & $0.705$ & $0.86$ & $0.602$ & $1.097$ & $0.543$ \\
deficit-gated, $\chi_r=0.05$ & $0.00$ & $0.00$ & $0.055$ & $1.334$ & $1.112$ & $1.136$ & $1.39$ & $1.112$ & $0.060$ & $0.145$ \\
deficit-gated, $\chi_r=0.10$ & $0.01$ & $0.01$ & $0.059$ & $1.331$ & $1.110$ & $1.135$ & $1.39$ & $1.110$ & $0.065$ & $0.139$ \\
deficit, $\chi_r=0.10$, conversion frozen & $0.05$ & $0.05$ & $0.088$ & $1.315$ & $1.096$ & $1.123$ & $1.37$ & $1.096$ & $0.088$ & $0.159$ \\
\bottomrule
\end{tabular}
\end{table}

Without restoration, $A_c$ stays at its minimum because the model has no return flow. Surplus-gated restoration produces partial or over-restoration depending on $\chi_r$. At $\chi_r\ge0.05$, it continues while $B>E$, and supported population falls to $0.97$--$0.90$ of its reference value. A target cap returns $A_c$ to its reference ($\mathcal R_c=1.00$). Deficit-gated restoration acts during shortfall, but at the tested rates it is weaker than concurrent conversion, leaving recovery near zero. Freezing conversion under surplus-gated restoration raises peak debt from $0.152$ to $0.543$\,gha\,yr; restoration then shifts reserve land to capital use while fast-land maintenance declines. These outcomes depend on the starting state and intervention rates. The gate-sign comparison keeps $T_b=0$; SI~\S S8 reports the separate matched wave/no-wave debt experiment.\par

\subsection{Regime table}\label{si:regimes}

\begin{table}[H]
\centering\small
\caption{The composition diagnostic $\dconv$ (gha\,ha$^{-1}$) as the decision condition for mask visibility. There is no interior area-based capacity maximum in the quality-based model; the only shape distinction is the sign of $\dconv$.}
\label{tab:si_regime}
\begin{tabularx}{\textwidth}{l l X}
\toprule
Regime & Condition & Behavior \\
\midrule
composition-premium & $\dconv>0$ & the conversion contribution to accounting $B$ is positive; net $B$ also depends on other flows (the mask is possible in the baseline) \\
composition-neutrality & $\dconv=0$ & first-order neutrality; non-generic \\
composition-discount & $\dconv<0$ & conversion lowers $B$ at the conversion margin; conversion alone cannot produce a composition mask \\
\bottomrule
\end{tabularx}
\end{table}

\section{Empirical details}\label{si:empirical}

\subsection{Data and vintages}\label{si:vintages}

Our empirical source is the National Ecological Footprint and Biocapacity Accounts (NFA), Edition 2026, results and metadata version 1.0 \citep{Lo2026NFA}. The workbook spans 1961--2025, but we use complete world observations through 2023 and exclude the estimated 2024--2025 rows. The 1961--2022 sensitivity uses the same edition and population series. This release recalculates historical values and reports five biocapacity account categories: fishing grounds, built-up land, cropland, grazing land, and forest biocapacity. These categories are account entries, not a common-boundary inventory of physical stocks.

The workbook's $\texttt{pop\_1000\_persons}$ field reports UN PopStat population in thousands; we convert it to persons before calculating per-capita values. The field $\texttt{efp\_total\_gha}$ gives the Global Ecological Footprint of Production, denoted $E_{\mathrm{prod}}$. We calculate per-capita values from unrounded account totals and population rather than the workbook's rounded columns. Table~\ref{tab:si_nfa_endpoints} summarizes the endpoints. From 1961 to 2023, $B$ rises from $9{,}725{,}417{,}562$ to $12{,}142{,}843{,}961$\,gha, population from $3.06487$ to $8.09173$ billion, and $B/P$ falls from $3.17319$ to $1.50065$\,gha\,person$^{-1}$. Over the same period, $E_{\mathrm{prod}}/P$ rises from $2.19700$ to $2.56462$\,gha\,person$^{-1}$. The identity $\Delta\ln B=\Delta\ln P+\Delta\ln(B/P)$ holds to rounding; it separates extensive and intensive changes but does not attribute causes. In both windows, 1973 is the first year with $E_{\mathrm{prod}}/B>1$.

\begin{table}[htbp]
\centering\small
\caption{World endpoint changes in NFA Edition 2026. The 1961--2022 column is a same-edition sensitivity; 2024--2025 estimates are excluded. Population and per-capita quantities are recomputed from workbook totals.}
\label{tab:si_nfa_endpoints}
\begin{tabularx}{\textwidth}{>{\raggedright\arraybackslash}Xrr}
\toprule
Measure & 1961--2023 & 1961--2022 \\
\midrule
$\Delta B/B_{1961}$ & $+24.8568\%$ & $+24.5450\%$ \\
$\Delta\ln B$ & $+0.221997$ & $+0.219497$ \\
$\Delta P/P_{1961}$ & $+164.0154\%$ & $+161.7211\%$ \\
$\Delta\ln P$ & $+0.970837$ & $+0.962109$ \\
$\Delta(B/P)/(B/P)_{1961}$ & $-52.7085\%$ & $-52.4131\%$ \\
$\Delta\ln(B/P)$ & $-0.748840$ & $-0.742612$ \\
$\Delta(E_{\mathrm{prod}}/P)/(E_{\mathrm{prod}}/P)_{1961}$ & $+16.7327\%$ & $+17.6994\%$ \\
First year $E_{\mathrm{prod}}/B>1$ & 1973 & 1973 \\
\bottomrule
\end{tabularx}
\end{table}

\newpage
\subsection{Divisia decomposition, sensitivity to the base-year share and level fallback}\label{si:divisia}

The crop proxy is an illustrative account construction. We combine independently sourced OWID/FAOSTAT world cropland area, defined as arable plus permanent-crop area (ha), with cereal yield (tonnes per hectare) \citep{OWIDCropland2026,OWIDCereal2026}. Cereal yield is a productivity sentinel, not a measure of output from all crops. For a declared 1961 account-share anchor $\phi_0$, we set $\hat b_f(t)=\hat b_f(1961)\tilde b_f(t)/\tilde b_f(1961)$, with $\hat b_f(1961)A_f(1961)=\phi_0B(1961)$. We then define $X_t=\hat b_f(t)A_f(t)$ and $C_t=B_t-X_t$. Thus $X$ is an index-scaled account proxy, and $C$ is the residual that reconciles it to the NFA total. The residual is not observed ecological capital, non-cropland biocapacity, or the model quantity $\Ups(q)A_c$.

When both $X$ and $C$ are positive at the two endpoints, the exact log-mean Divisia identity is
\begin{equation}
\Delta\ln B=w_X\bigl(\Delta\ln\tilde b_f+\Delta\ln A_f\bigr)+w_C\Delta\ln C,
\qquad w_X=\frac{L(X_T,X_0)}{L(B_T,B_0)},\quad w_C=\frac{L(C_T,C_0)}{L(B_T,B_0)},
\end{equation}
Here $L(a,b)=(b-a)/(\ln b-\ln a)$ is the logarithmic mean, with its continuous limit when $a=b$. We use the measured 1961 NFA cropland share, $\phi_0=0.1286472202$, as the central anchor. Mapping that account share to the external area--yield proxy is an assumption, not an independent observation of physical capital. Table~\ref{tab:si_phi0} gives the selected sensitivity for 1961--2023. At the central anchor, the cereal-yield, area, and residual terms are $+0.268662$, $+0.037710$, and $-0.084375$, summing to $\Delta\ln B=+0.221997$. The same-edition 1961--2022 sensitivity gives $+0.264696$, $+0.037506$, and $-0.082705$, summing to $+0.219497$. The decomposition is exact for this proxy, but its allocation depends on $\phi_0$ and is not a physical-stock attribution.

\begin{table}[htbp]
\centering\footnotesize
\caption{Sensitivity of the proxy Divisia decomposition to the 1961 account-share anchor $\phi_0$ (world, 1961--2023). Contributions and weights are dimensionless. The $\phi_0=0.128647$ row uses the measured NFA cropland share in 1961; other rows are illustrative calibrations.}
\label{tab:si_phi0}
\begin{tabular}{r r r r r r r}
\toprule
$\phi_0$ & $w_X$ & $w_C$ & Cereal yield & Area & Residual $C$ & Sum \\
\midrule
$0.100000$ & $0.1833$ & $0.7957$ & $+0.2088$ & $+0.0293$ & $-0.0162$ & $+0.2220$ \\
$0.128647$ & $0.2358$ & $0.7352$ & $+0.2687$ & $+0.0377$ & $-0.0844$ & $+0.2220$ \\
$0.150000$ & $0.2749$ & $0.6893$ & $+0.3133$ & $+0.0440$ & $-0.1352$ & $+0.2220$ \\
$0.190000$ & $0.3483$ & $0.6008$ & $+0.3968$ & $+0.0557$ & $-0.2305$ & $+0.2220$ \\
$0.250000$ & $0.4582$ & $0.4581$ & $+0.5221$ & $+0.0733$ & $-0.3734$ & $+0.2220$ \\
$0.300000$ & $0.5499$ & $0.3178$ & $+0.6265$ & $+0.0879$ & $-0.4924$ & $+0.2220$ \\
$0.320000$ & $0.5866$ & $0.2454$ & $+0.6683$ & $+0.0938$ & $-0.5401$ & $+0.2220$ \\
$0.340000$ & $0.6232$ & $0.1009$ & $+0.7100$ & $+0.0997$ & $-0.5877$ & $+0.2220$ \\
\bottomrule
\end{tabular}
\end{table}

For $\phi_0\ge0.340531$ in 2023 (or $0.343400$ in 2022), terminal residual $C$ is nonpositive, so logarithmic weighting is invalid. We use the exact level identity $\Delta B=\Delta X+\Delta C$ instead; this is still an accounting identity, not a stock measure. At the selected anchor, $X_{2023}/B_{2023}=0.377784$ and $C$ remains positive, but the proxy fast-land share and residual contribution vary across the sensitivity analysis. Every row uses the same NFA total $B$. The residual absorbs differences between the index-scaled crop proxy and that total, including account categories, boundaries, valuation, vintage, and calibration.

\subsection{Cropland-account calibration}\label{si:cropland}

We report NFA cropland-account movements separately from the illustrative proxy. From 1961 to 2023, NFA cropland biocapacity rises $2.870055$-fold and NFA cropland area rises $1.198029$-fold; their account-value-per-area quotient rises $2.395647$-fold. For 1961--2022, the same-edition ratios are $2.851441$, $1.198326$, and $2.379520$. This quotient is an NFA account value per area, not an independent physical crop-yield measure. Over 1961--2023, the external OWID/FAOSTAT area series rises $1.173414$-fold and cereal yield rises $3.124667$-fold. Those independent series serve as a sentinel proxy; they do not reconstruct the NFA cropland account or measure output from all crops. The model's representative $b_{f0}=0.85\,\mathrm{gha\,ha^{-1}}$ remains a conditional parameter, not an empirical fit to these ratios.

\subsection{Observation-operator hierarchy}\label{si:operator}

\begin{itemize}[leftmargin=1.5em]
\item \emph{NFA accounting observations.} Total and component values in gha describe NFA account composition and change under the stated edition and conventions. They do not directly measure biomass, soil condition, wetland extent, fish biomass, or model state $A_c$.
\item \emph{Fast-provisioning area.} The external world cropland-area series covers arable land plus permanent crops for 1961--2023 (ha) \citep{OWIDCropland2026}. It is an independent proxy input, not the NFA cropland-area series.
\item \emph{Crop-productivity sentinel.} OWID/FAOSTAT cereal yield is reported in tonnes per hectare for 1961--2024 \citep{OWIDCereal2026}. It covers cereals, not all crops, and is not a physical-stock measure. Base-anchored $\hat b_f$ is used only in the illustrative account proxy.
\item \emph{Ecological-capital components.} Separate physical observation operators could use forest extent and biomass, soil carbon and condition, wetland extent and condition, and fish biomass by species or stock. Such data would need harmonized boundaries, reference years, definitions, uncertainty, and explicit treatment of overlap. We assembled no common-boundary global physical-stock series for these components. NFA forest biocapacity and residual $C$ cannot substitute for those measurements.
\item \emph{Model mapping.} The model defines capital contribution as $Y_c=\Ups(q)A_c$. Neither $A_c$ nor $q$ is identified empirically by total $B$ or an NFA component, and we do not assume proxy residual $C=B-X$ equals $Y_c$. A model--data comparison would first need independent physical measures and an explicit mapping.
\end{itemize}

Table~\ref{tab:si_model_nfa} summarizes the possible links while keeping NFA categories distinct from model states.

\begin{table}[htbp]
\centering\small
\caption{Conceptual mapping from two-land model quantities to NFA information and independent proxy inputs. The rows specify possible observation operators, not identities between NFA categories and model states. Annual account values are in gha; land areas are in ha.}
\label{tab:si_model_nfa}
\begin{tabularx}{\textwidth}{>{\raggedright\arraybackslash}p{0.23\textwidth} >{\raggedright\arraybackslash}p{0.31\textwidth} X}
\toprule
Model quantity & Closest account or proxy information & Interpretation and identification limit \\
\midrule
Fast area and proxy account $A_f$, $X=\hat b_fA_f$ & NFA cropland account; independently sourced OWID cropland area and cereal yield & NFA cropland biocapacity is an account value; the external area--yield product is an illustrative, base-anchored proxy. Neither is an observed model fast-land yield or a one-to-one physical measure. \\
Capital area, quality and contribution $A_c$, $q$, $Y_c=\Ups(q)A_c$ & No direct NFA counterpart; the proxy residual $C=B-X$ is defined by subtraction & In the model $Y_c=\Ups(q)A_c$; empirically $C$ is not assumed equal to $Y_c$. It also absorbs differences in boundaries, account categories (including built-up land), valuation, vintage and proxy calibration, and is not a physical-stock measure. \\
Reserve area $A_r$ & No direct NFA account counterpart & $A_r$ closes the model's conserved-area ledger; it is not recoverable from the NFA aggregate. \\
\bottomrule
\end{tabularx}
\end{table}

\begin{table}[htbp]
\centering\small
\caption{Observability ladder: what is directly measured, what is identifiable by construction, and what remains unidentified.}
\label{tab:si_ladder}
\begin{tabularx}{\textwidth}{>{\raggedright\arraybackslash}p{0.3\textwidth} X >{\raggedright\arraybackslash}p{0.3\textwidth}}
\toprule
Tier & Quantities & Status \\
\midrule
Observable physical area input & OWID cropland area [ha]; other mapped land-cover classes where available & Measured extent, but mapping to effective $A_c$ is not direct \\
Observable population & $P$ [persons] & Direct in the UN PopStat workbook field after unit conversion \\
Observable productivity sentinel & Cereal yield [tonnes ha$^{-1}$] & Physical yield of cereals only; not all-crops output or ecological capital \\
Observable annual account value & NFA $B$, $E_{\mathrm{prod}}$ and components [gha] & Measured only under the NFA accounting convention; value-weighted accounts \\
Proxy account and residual & $X=\hat b_fA_f$ and $C=B-X$ [gha] & Identified by construction for the stated inputs and anchor; not physical stocks and not assumed equal to $\Ups(q)A_c$ \\
Effective capital contribution & $Y_c=\Ups(q)A_c$ [gha] & A model quantity; no common-boundary physical observation operator is assembled here \\
Not separately identifiable & $A_c$, $q$, $b_c$, $b_{G,c}$, $\rho_c$ and $D$ & Requires independent physical-state, quality, regeneration or degradation observations and a declared mapping \\
\bottomrule
\end{tabularx}
\end{table}

\subsection{Regeneration anchors}\label{si:anchors}

The model's capital-land regeneration timescale represents a lumped recruitment response, not a universal ecological rate. In Neotropical secondary forests, plot-level age and above-ground biomass records (2ndFOR; 1,334 plots, 41 chronosequences, 35 good fits) give median $t_{50}\approx30$\,yr and $\tau_r\approx38$\,yr under a lagged saturating recovery model \citep{Poorter2016}. For 33 fish stocks classified as collapsed in RAM Legacy, trajectories of $B/B_{MSY}$ give $t_{50}\approx5.6$\,yr and $\tau_r\approx3.4$\,yr \citep{HutchingsReynolds2004,Neubauer2013}. Soil organic carbon re-equilibrates over about 17--23\,yr after land-use change \citep{Poeplau2011}. Forest biomass reaches about half of old-growth levels by 20\,yr and 90\% at a median of about 66\,yr; roughly 29\% of collapsed stocks recover to 50\% within 5--15\,yr.

These observations describe different systems and are not pooled into one rate. They are examples for future ecosystem-specific measurement and modeling, not a basis for assigning a generic response time to the lumped capital-land class. The 10--40\,yr interval is a scenario range, not a pooled empirical estimate; the faster fish-stock episodes are not treated as a lower bound for the generic class. Neither $\rho_c$ nor $\tau_g$ is calibrated from these records, and they do not validate the present quality equation or parameter values.

\subsection{Recursive identifiability in the land-type accounts}\label{si:recursive}

NFA Edition 2026 illustrates why the aggregate should be read with its components, without treating account movements as physical-stock changes. From 1961 to 2023, cropland contributes $+24.0577$ percentage points to the $+24.8568\%$ change in total $B$; built-up land adds $+7.4416$, while Forest biocapacity contributes $-5.9525$, grazing land $-0.2688$, and fishing grounds $-0.4212$. The non-cropland sum, which includes built-up land, changes by only $+0.7990$ points ($1.0092\times$) because unlike account movements offset one another.

For the same-edition 1961--2022 sensitivity, contributions are $+23.8183$ points for cropland, $+7.2854$ for built-up land, $-5.8137$ for Forest biocapacity, $-0.2890$ for grazing land, and $-0.4559$ for fishing grounds. The non-cropland sum changes by $+0.7267$ points ($1.0083\times$). These are accounting attributions, not evidence that forest biomass, soil condition, wetland extent, fish biomass, or model $A_c$ declined. Built-up land is an NFA account class, not ecological capital; including it in the non-cropland sum is arithmetic, not an ecological aggregation.

\subsection{Physical-stock scoreboard: specification, not inventory}\label{si:stocks}

We assembled no common-boundary global physical-stock series for these components, so we present no physical-stock composite or global decline finding. This does not mean that such a series could never be assembled. Table~\ref{tab:si_stock_scoreboard} is an unpopulated specification of candidate measures and comparability requirements; it contains no empirical endpoints and is not a model calibration.

These physical measures are heterogeneous. Forest biomass and soil carbon are masses; wetland and cropland extent are areas; fish biomass is stock-specific; freshwater storage is a volume; and condition indicators are often dimensionless or categorical. They are not dimensionally comparable and may overlap spatially. Annual flows—harvest, catch, increment, recharge, or emissions—must also remain distinct from stocks and extents. Even after normalization, interpretation depends on overlap, baselines, uncertainty, ecological thresholds, valuation, and allowed substitution. Unless common spatial and temporal boundaries, a normalization or valuation rule, and substitution rules are specified, the components should remain a nonadditive vector. A scalar composite is not neutral and cannot be inferred from NFA gha accounts. SEEA Ecosystem Accounting likewise distinguishes ecosystem extent and condition from ecosystem-service and monetary presentations \citep{UN2021SEEA_EA}.

\begin{table}[htbp]
\centering\footnotesize
\caption{Physical-stock scoreboard specification. The rows are unpopulated: no common-boundary global physical-stock series for the relevant components or endpoint values were assembled. Stocks, extents, condition measures, and annual flows are not combined into a composite.}
\label{tab:si_stock_scoreboard}
\begin{tabularx}{\textwidth}{@{}>{\raggedright\arraybackslash}p{0.12\textwidth} >{\raggedright\arraybackslash}p{0.27\textwidth} >{\raggedright\arraybackslash}p{0.22\textwidth} >{\raggedright\arraybackslash}X@{}}
\toprule
Component & Physical stock, extent, or condition to define & Annual flow or pressure (separate) & Comparability rule required \\
\midrule
Forests & Extent (ha); growing stock or biomass (m$^3$ or tonnes C); stated condition & Net annual increment, harvest, or deforestation & Common forest definition and mapped area; reference year; treatment of plantations and overlap \\
Soils & Soil organic-carbon stock or density for a stated depth; structure or erosion condition & Carbon accumulation or loss; erosion where measured & Consistent depth, bulk density, land cover, and sampling method \\
Wetlands & Extent (ha) and an explicit ecological-condition measure & Carbon accumulation, water-storage change, or nutrient flux & Wetland definition and overlap with forest and other land-cover classes \\
Fisheries & Spawning-stock biomass by species or stock (tonnes); age or size structure & Recruitment, catch, and landings & Stock identity, assessment method, reference point, and uncertainty \\
Freshwater & Surface-water or groundwater storage (m$^3$), level, and quality & Recharge and withdrawal & Basin or aquifer boundary; separate stored volume from annual flow \\
\bottomrule
\end{tabularx}
\end{table}

In the two-land model, $A_c$ is an effective area state (ha) and $q$ is a separate quality index; neither is observed in the NFA total. Physical-stock data could inform a model observation operator only after boundaries and mappings to these effective states are specified. Proxy residual $C=B-X=B-\hat b_fA_f$ is not ecological capital, a measured non-cropland stock, or an estimate of $\Ups(q)A_c$. NFA components remain account values, not physical-stock observations.
\section{Policy extensions}\label{si:policy}

\paragraph{Reserved biocapacity share.} If both land classes have the same reservation share, $\sigma_f=\sigma_c=\sigma<1$, human-available supply is $\sigma B$ and shortfall is $S=[E-\sigma B]_+$. Reservation lowers the demand at which conversion begins; conversion is off while $E\le\sigma B$. With different shares, the condition becomes $E\le\sigma_fY_f+\sigma_cY_c$, which need not equal $\sigma B$. These results apply to the stated parameterization; they do not establish robustness to disturbances or parameter uncertainty.

\paragraph{Typed floors and conversion stop.} The model imposes separate constraints on reserve land ($A_r\ge A_r^{\min}$), capital land ($A_c\ge\Acmin$), fast land ($A_f\le A_f^{\max}$), and carrying capacity ($K\ge K_{\min}$). A stop rule triggered at $A_c=\Acmin$ preserves the capital-land floor. A ceiling on $A_f$ limits fast-land extent but does not ensure that $A_c$ stays above its floor unless the stop rule also accounts for the capital floor and remaining-area constraints. In Eq.~(5), conversion slows to zero as $A_c\to\Acmin$. A control based only on total $B$ cannot guarantee this typed floor because $B$ does not identify $A_c$.

\paragraph{Institutional signature.} Related work distinguishes protective and mobilizing institutional feedback \citep{Abaee2026a}. Here, surplus-gated restoration is off during shortfall, while deficit-gated restoration acts during shortfall and can further reduce provision (Section~\ref{si:gates}).

\paragraph{Aggregation and typed feasibility.} A related transition-assessment study defines a linear compensatory operator and compares perfect substitutability with less-substitutable CES alternatives \citep[Sections~2 and 4 of][]{Abaee2026b}. ECOMOD's $B$ is also a linear sum, but the companion considers a separate finite-plan setting. There, the compensatory protocol accepts states even when no available plan meets every component floor (Theorem~2 in \citet{Abaee2026b}). Theorem~4(iii) in \citet{Abaee2026b} identifies an exact interval of diagonal-gap states on its witness-defined region that each positive elasticity falsely certifies as safe. Thus no positive-elasticity member is uniformly safe on that region. Under pathwise feasibility, falling below a typed floor during transition remains a violation even if the endpoint later recovers. The comparison distinguishes an aggregate criterion from constraints on individual components.

\section{The one-stock comparator}\label{si:onestock}

We include this one-stock formulation as a structural comparator, not as a reduction of the two-land model. Its state $A$ has comparator-specific stock units and is distinct from two-land state $A_c$. The equations are
\begin{align*}
G(A)&=\rho A\left(1-\frac{A}{A_{\max}}\right),\qquad
B(t)=b(t)A(t)+b_GG(A(t)),\qquad E(t)=eP(t);\\
\dot A(t)&=G\bigl(A(t-\tau_g)\bigr)-\frac{\pos{E(t)-\sigma b(t)A(t)}}{b_G};\\
\dot P(t)&=rP(t)\left(1-\frac{P(t)}{K(t-\tau_p)}\right),\qquad
K(t)=\max\!\left(\frac{B(t)}{e},K_{\min}\right);\\
\dot D(t)&=\pos{E(t)-B(t)}-\eta D(t),\qquad
b(t)=\bigl(b_0+T_b(t)\bigr)e^{-\alpha D(t)};\\
T_b(t)&=\omega_{Tb}\frac{\Delta b}{1+\exp[-\kappa(t-t_w)]}.
\end{align*}
Here $\sigma$ is the dimensionless human-available share; $\omega_{Tb}\in\{0,1\}$ switches the logistic technology increment on or off. The states $A$, $A_{\max}$, and imposed floor $A_{\mathrm{ext}}$ use comparator-specific stock units. $A_{\max}$ sets the logistic growth scale; it is not an imposed upper bound. Accordingly, $G(A)$ is in stock units per year. The coefficients $b$, $b_0$, $T_b$, and $\Delta b$ are in gha per stock unit, while $b_G$ is in gha\,yr per stock unit, making both terms in $B$ annual account values in gha.

$B$ and $E$ use the GFN/NFA convention, with the annual basis implicit. Population $P$, carrying capacity $K$, and $K_{\min}$ are in cap; $e$ is in gha\,cap$^{-1}$. Debt $D$ is in gha\,yr, so $\dot D$ and $\eta D$ are in gha, and $\alpha$ is in gha$^{-1}$\,yr$^{-1}$. Rates $\rho,r,\eta,\kappa$ are in yr$^{-1}$; delays $\tau_g,\tau_p$ and wave midpoint $t_w$ are in years.

The numerical implementation enforces $A\ge A_{\mathrm{ext}}$, $P,D\ge0$, and $K\ge K_{\min}$, with the exact positive-part operator. Delayed histories are constant at their initial values, and RK4 integration uses linear interpolation of delayed states. Step sizes are $0.1$\,yr for the full trajectory (checked at $0.05$ and $0.025$\,yr), $0.01$\,yr for the fixed-demand mask (with a $0.02$\,yr shortfall sweep), $0.2$\,yr for the delay pair (checked at $0.1$\,yr), and $0.5$\,yr for the finite-grid classification.

With fixed $b$, $T_b=0$, $D=0$, and $\sigma=1$, interior equilibria satisfying $0<A<A_{\max}$ and $B(A)/e>K_{\min}$ obey $E=B(A)$ and $P=B(A)/e$. They form a one-parameter family along
\[
B(A)=(b+b_G\rho)A-\frac{b_G\rho}{A_{\max}}A^2,\qquad
B'(A)=b+b_G\rho-2b_G\rho\frac{A}{A_{\max}}.
\]
An interior maximum within $0<A<A_{\max}$ requires $b_G\rho>b$. For the debt-free, constant-parameter delay and basin illustrations, $b=0.5$, $b_G=0.8$, and $\rho=0.05$, so $b_G\rho=0.04<b$ and $B(A)$ increases throughout that interval. In the full-trajectory illustration, $b(t)$ changes with the technology wave and debt; we make no monotonicity or equilibrium-fold claim for that time-varying case.

For the debt-free reference equilibrium, $A^*=0.8$ under $(b,b_G,\rho,A_{\max},e,r)=(0.5,0.8,0.05,1.2,0.55,0.02)$. It has $B^*=0.4107$\,gha and $P^*=B^*/e=0.7467$\,cap. The characteristic equation has a neutral root tangent to the equilibrium family and a positive real root. That positive root is $0.5883$\,yr$^{-1}$ when $(\tau_g,\tau_p)=(0,0)$ and $0.6250$\,yr$^{-1}$ when $(30,25)$\,yr. These local results apply only to the comparator; they are neither a Hopf classification nor a stability result for the two-land system.

Figures~\ref{fig:si_onestock_full}--\ref{fig:si_onestock_basin} show conditional runs of the comparator. In Fig.~\ref{fig:si_onestock_delay}, the same starting state approaches $A_{\max}$ at $\tau_g=10$\,yr but reaches the imposed floor at $\tau_g=30$\,yr, with $\tau_p=25$\,yr in both cases. This pair shows delay sensitivity for one initial condition; it does not establish a universal critical delay. Figure~\ref{fig:si_onestock_basin} classifies a finite grid of starting states at fixed $\tau_p$. Grid-cell proportions are not probabilities over initial conditions or measures of a continuous basin.

\begin{figure}[htbp]
\centering
\includegraphics[width=0.96\textwidth]{si_onestock_full_trajectory}
\caption{Full delayed one-stock comparator with a logistic technology increment. Parameters are $A_0=1.0$ stock unit, $P_0=0.45$\,cap, $b_0=0.5$, $b_G=0.6$, $\rho=0.08$\,yr$^{-1}$, $A_{\max}=1.2$ stock units, $e=0.55$\,gha\,cap$^{-1}$, $r=0.02$\,yr$^{-1}$, $\eta=0.03$\,yr$^{-1}$, $\alpha=0.2$\,gha$^{-1}$\,yr$^{-1}$, $\tau_g=15$\,yr, $\tau_p=25$\,yr, $A_{\mathrm{ext}}=0.02$ stock units, $K_{\min}=10^{-4}$\,cap, $\Delta b=0.5$\,gha per stock unit, $\kappa=0.04$\,yr$^{-1}$, and $t_w=120$\,yr; $\sigma=1$, $\omega_{Tb}=1$, and $D_0=0$. The positive-part operator is exact. $B$ and $E$ are in gha, $A$ is in comparator-specific stock units, $P$ is in cap, and $D$ is in gha\,yr. The trajectory reaches the imposed lower bound near $721$\,yr, then recovers as modeled demand falls; the bound is not absorbing. This is a conditional model illustration, not an observation or a two-land result.}
\label{fig:si_onestock_full}
\end{figure}

\begin{figure}[htbp]
\centering
\includegraphics[width=0.96\textwidth]{si_onestock_mask_window}
\caption{One-stock masking example with fixed demand and no delays. Parameters are $A_0=1.0$ stock unit, $E=0.60$\,gha, $b_0=0.5$, $b_G=0.8$, $\rho=0.05$\,yr$^{-1}$, $A_{\max}=1.2$ stock units, $\eta=0.05$\,yr$^{-1}$, $\alpha=0.03$\,gha$^{-1}$\,yr$^{-1}$, $\Delta b=1.5$\,gha per stock unit, $\kappa=0.2$\,yr$^{-1}$, and $t_w=15$\,yr; $\sigma=1$, $\omega_{Tb}=1$, and $D_0=0$. With the exact positive-part operator, the interval with $\dot B>0$ and $\dot A<0$ lasts about $3.0$\,yr. Over that interval, $A$ falls from about $0.937$ to $0.867$ stock units while $B$ rises from about $0.575$ to $0.600$\,gha. Panel (b) shows the interval width across initial shortfalls $[E(0)-b(0)A_0]_+$ for this parameterization. These are conditional comparator outputs.}
\label{fig:si_onestock_mask}
\end{figure}

\begin{figure}[htbp]
\centering
\includegraphics[width=0.96\textwidth]{si_onestock_delay_examples}
\caption{Controlled regeneration-delay examples in the debt-free one-stock subsystem. Both runs use $A_0=1.0$ stock unit, $P_0=0.1$\,cap, $b=0.5$, $b_G=0.8$, $\rho=0.05$\,yr$^{-1}$, $A_{\max}=1.2$ stock units, $e=0.55$\,gha\,cap$^{-1}$, $r=0.02$\,yr$^{-1}$, $\tau_p=25$\,yr, $A_{\mathrm{ext}}=0.02$ stock units, $K_{\min}=10^{-6}$\,cap, and $\sigma=1$; only $\tau_g$ differs ($10$ versus $30$\,yr). $B$ is in gha, $P$ in cap, and $A$ in comparator-specific stock units. These selected trajectories do not establish a general critical delay.}
\label{fig:si_onestock_delay}
\end{figure}

\begin{figure}[htbp]
\centering
\includegraphics[width=0.96\textwidth]{si_onestock_basin}
\caption{Finite-grid initial-condition classification in the debt-free one-stock subsystem. The runs use the parameter set in Fig.~\ref{fig:si_onestock_delay}, with $K_{\min}=10^{-6}$\,cap and $\tau_p=25$\,yr; $\tau_g=0$\,yr in the left panel and $30$\,yr in the right. The imposed lower bound is $A_{\mathrm{ext}}=0.02$ stock units. The grid uses $A_0=0.10,0.20,\ldots,1.30$ stock units and $P_0=0.05,0.15,\ldots,1.55$\,cap; trajectories run to $1200$\,yr with $\Delta t=0.5$\,yr. Green marks a high endpoint $A(1200)\ge0.95A_{\max}$, red a low endpoint $A(1200)\le0.05$ stock units, and gray an intermediate endpoint. These display-grid categories are not probabilities or continuous-basin measures.}
\label{fig:si_onestock_basin}
\end{figure}

\section{Controlled technology-wave experiment}\label{si:tbwave}

The matched experiment compares endogenous-demand runs with the technology increment $T_b$ switched on or off, at $P_0=1.0$ or $1.2$\,cap. Within each pair, the runs share $A_{tot}=2.5$\,ha, $A_{f0}=0.6$\,ha, $A_{c0}=0.9$\,ha, $q_0=0.9$, and $D_0=0.0$\,gha\,yr. Parameters and numerical methods are matched; only $\omega_{Tb}$ differs. Both runs also share constant pre-run histories at $q_0=0.9$ and the no-wave initial account value $B=0.560184$\,gha.

Shared parameter values are $b_{f0}=0.85$\,gha\,ha$^{-1}$, $b_c=0.05$\,gha\,ha$^{-1}$, $b_{G,c}=0.8$\,gha\,yr\,ha$^{-1}$, $\rho_c=0.08$\,yr$^{-1}$, $\alpha=0.03$\,gha$^{-1}$\,yr$^{-1}$, $\eta=0.05$\,yr$^{-1}$, $\eta_f=0.05$\,yr$^{-1}$, $\mu=0.06$\,yr$^{-1}$, $e=0.55$\,gha\,cap$^{-1}$, $r=0.02$\,yr$^{-1}$, $\tau_g=20$\,yr, $\tau_p=25$\,yr, $\tau_{conv}=1$\,yr, and $\chi_r=0.10$\,gha$^{-1}$\,yr$^{-1}$. The active wave uses $\Delta b=1.0$\,gha\,ha$^{-1}$, $\kappa_w=0.1$\,yr$^{-1}$, and $t_{wave}=40$\,yr; runs last $250$\,yr. The logistic increment is already nonzero at $t=0$ ($T_b(0)=0.017986$\,gha\,ha$^{-1}$), so initial $B$ is $0.560184$\,gha without the wave and $0.570976$\,gha with it. This small difference comes from the logistic tail. Land, quality, population, and debt states otherwise match, and the pre-run $B$ history remains common. The wave reaches half its increment at $t_{wave}=40$\,yr.

Figure~\ref{fig:si_wave_debt} shows the $P_0=1.0$ pair. In both runs, peak debt occurs after the wave midpoint: $D_{max}=0.008126$\,gha\,yr at $86.15$\,yr without the wave and $0.020205$\,gha\,yr at $115.60$\,yr with it. At $P_0=1.2$, the all-time maximum occurs early and is lower with the wave, but the maximum for $t\ge t_{wave}$ is higher. Table~\ref{tab:si_wave_debt} reports both windows. Because $t_{wave}$ marks the logistic midpoint, the post-midpoint window includes the wave's continued rise.

\begin{table}[htbp]
\centering\small
\caption{Matched technology-wave/no-wave debt maxima. Entries are $D_{max}$ in gha\,yr, with the time of the maximum in years in parentheses; the ratio is wave/no-wave. ``All'' is the maximum over $0\le t\le250$\,yr; $t\ge40$ is the post-midpoint maximum. The $P_0=1.0$ all-time peaks both occur after the midpoint.}
\label{tab:si_wave_debt}
\begin{tabularx}{\textwidth}{r l r l >{\centering\arraybackslash}X >{\centering\arraybackslash}X r}
\toprule
$P_0$ (cap) & quality update & $\Delta t$ (yr) & window & no wave $D_{max}$ (time) & wave $D_{max}$ (time) & ratio \\
\midrule
$1.0$ & delayed & $0.050$ & all & $0.008126\ (86.15)$ & $0.020205\ (115.60)$ & $2.487$ \\
$1.0$ & delayed & $0.025$ & all & $0.008105\ (86.20)$ & $0.020195\ (115.65)$ & $2.492$ \\
$1.0$ & current & $0.050$ & all & $0.008326\ (85.70)$ & $0.020168\ (115.60)$ & $2.422$ \\
$1.2$ & delayed & $0.050$ & all & $0.061218\ (1.65)$ & $0.051037\ (1.45)$ & $0.834$ \\
$1.2$ & delayed & $0.050$ & $t\ge40$ & $0.010854\ (81.70)$ & $0.025184\ (112.55)$ & $2.320$ \\
$1.2$ & delayed & $0.025$ & all & $0.060732\ (1.625)$ & $0.050582\ (1.475)$ & $0.833$ \\
$1.2$ & delayed & $0.025$ & $t\ge40$ & $0.010808\ (81.775)$ & $0.025165\ (112.625)$ & $2.328$ \\
\bottomrule
\end{tabularx}
\end{table}

For $P_0=1.0$, the all-time wave/no-wave peak-debt ratio is $2.487$ at $\Delta t=0.05$\,yr and $2.492$ at $0.025$\,yr; under current-quality regeneration it is $2.422$ at $0.05$\,yr. For $P_0=1.2$, all-time ratios are $0.834$ and $0.833$, while post-midpoint ratios are $2.320$ and $2.328$. Halving the step changes debt maxima by less than $0.9\%$. In the $P_0=1.0$ pair, $A_c$ stays at or above its initial $0.9$\,ha, so this experiment illustrates a conditional delayed debt rebound, not the frozen-demand capital-land drawdown. At $P_0=1.2$, the all-time and post-midpoint maxima tell different stories. The debt variable $D$ is a leaky state driven by aggregate deficit $E-B$, not an observed NFA or physical-stock quantity.

\begin{figure}[!htbp]
\centering
\includegraphics[width=0.92\textwidth]{fig_tb_wave_debt}
\caption{Matched endogenous-demand technology-wave/no-wave simulation at $P_0=1.0$\,cap. Parameters and initial conditions are in SI~\S S8; only $T_b$ is switched on or off. (a) The active logistic increment $T_b(t)$, with the dotted line at its midpoint $t_{wave}=40$\,yr; (b) annual account values $B$ and $E$ (gha); (c) population $P$ (cap); and (d) leaky model debt $D$ (gha\,yr). Markers denote all-time maxima: $86.15$\,yr without the wave and $115.60$\,yr with it; peak debt is $0.008126$ versus $0.020205$\,gha\,yr. The wave begins with its small logistic tail at $t=0$ and approaches $\Delta b=1.0$\,gha\,ha$^{-1}$. The dotted line marks the midpoint, not an end date.}
\label{fig:si_wave_debt}
\end{figure}

